\chapter{AI Klasik: Mencari, Merencanakan, Memutuskan}
\label{ch:klasik}

Kuliah pengantar AI baru saya ikuti di semester terakhir, pada musim panas 2026, setelah hampir
semua kuliah deep learning selesai. Dengan urutan terbalik seperti itu, search, planning, game
playing, dan representasi pengetahuan tidak tampil sebagai masa lalu yang usang, melainkan sebagai
fondasi yang masih dipakai di dalam sistem modern, dari AlphaGo sampai agen berbasis model bahasa. Setahun sebelumnya, kuliah operations
research memberi sisi lain dari cerita yang sama: mencari keputusan terbaik di antara pilihan
yang terhingga tetapi banyak sekali.

Sebelum jaringan saraf mendominasi, AI berarti program yang bernalar secara eksplisit: mencari
jalan di ruang kemungkinan, menyusun rencana langkah demi langkah, dan menarik kesimpulan dari
fakta. Bab ini membahas gagasan-gagasan itu, karena tanpa mereka sulit memahami mengapa sistem AI
paling kuat hari ini menggabungkan pembelajaran dengan pencarian.

\section{Dari Turing ke Dartmouth}

\citet{turing1950} mengganti pertanyaan ``bisakah mesin berpikir?'' dengan permainan tiruan:
kalau seorang penanya tidak bisa membedakan mesin dari manusia lewat percakapan tertulis, apa
alasan kita menolak menyebut mesin itu cerdas? Istilah \emph{artificial intelligence} sendiri
lahir dari proposal lokakarya musim panas di Dartmouth pada 1956. Sejak itu, sejarah AI bergerak
dalam gelombang optimisme dan kekecewaan. Sistem pakar pada 1980-an menyimpan pengetahuan
manusia dalam ribuan aturan jika-maka, lalu tersandung pada sulitnya menulis semua aturan itu
dengan tangan. Gelombang statistik dan kemudian deep learning menggeser fokus dari aturan yang
ditulis ke pola yang dipelajari \citep{russell2021}.

Kuliah juga menyinggung perdebatan filosofisnya. \citet{searle1980} mengajukan eksperimen pikiran
ruang Cina. Seseorang yang tidak mengerti bahasa Cina duduk di dalam ruangan dan menjawab pesan
dalam bahasa Cina dengan mengikuti buku aturan yang lengkap. Dari luar, jawabannya
sempurna. Dari dalam, tidak ada pemahaman sama sekali. Searle berargumen bahwa memanipulasi simbol
tidak sama dengan memahami. Perdebatan antara AI lemah, yang hanya meniru perilaku cerdas, dan AI
kuat, yang sungguh memahami, belum selesai sampai sekarang, dan model bahasa besar membuatnya
kembali hangat.

\section{Pencarian}

Banyak masalah bisa dirumuskan sebagai pencarian: ada keadaan awal, sekumpulan aksi, fungsi
transisi yang mengatakan ke mana sebuah aksi membawa kita, uji tujuan, dan biaya setiap langkah.
Teka-teki geser delapan angka, rute navigasi, dan penyusunan jadwal semuanya cocok dengan bentuk
ini. Solusinya adalah jalur dari keadaan awal ke keadaan tujuan.

Pencarian tanpa informasi tidak tahu apa-apa tentang di mana tujuan berada. \istilah{Breadth-first
search} menjelajah lapis demi lapis dan menjamin jalur terpendek dalam jumlah langkah, tetapi
memorinya tumbuh eksponensial. \istilah{Depth-first search} masuk sedalam mungkin lebih dulu,
hemat memori, tetapi bisa tersesat di cabang yang tak berujung. \istilah{Uniform-cost search}
selalu memperluas simpul dengan biaya kumulatif terkecil, sehingga menemukan jalur termurah.

Pencarian dengan informasi memakai \istilah{heuristic} $h(n)$, perkiraan biaya dari simpul $n$ ke
tujuan. Algoritme A* \citep{hart1968} memperluas simpul dengan nilai
\begin{equation}
f(n)=g(n)+h(n)
\end{equation}
terkecil, dengan $g(n)$ biaya yang sudah ditempuh. Kalau heuristik tidak pernah melebihi biaya
sebenarnya, yaitu \istilah{admissible}, A* dijamin menemukan jalur optimal. Pada teka-teki delapan
angka, jumlah ubin yang salah tempat adalah heuristik admissible, dan jumlah jarak Manhattan setiap
ubin ke posisi seharusnya juga admissible tetapi lebih informatif, karena selalu sama atau lebih
besar. Heuristik yang lebih informatif memangkas jauh lebih banyak simpul.

A* mirip pengemudi yang memilih belokan berdasarkan dua hal sekaligus: jarak yang sudah ditempuh
di odometer, dan jarak garis lurus yang tersisa ke tujuan. Karena garis lurus tidak pernah lebih
panjang dari jalan sebenarnya, ia tidak akan tertipu untuk mengabaikan rute yang sebenarnya baik.

Ketika jalurnya tidak penting dan yang dicari hanya keadaan terbaik, misalnya susunan jadwal yang
paling sedikit bentrok, \istilah{local search} lebih cocok. Hill climbing terus bergerak ke tetangga
yang lebih baik sampai tidak ada lagi, dan sering terjebak di puncak lokal. Simulated annealing
\citep{kirkpatrick1983} kadang menerima langkah yang memburuk dengan peluang yang menurun perlahan,
seperti logam yang didinginkan pelan-pelan agar atomnya menemukan susunan berenergi rendah. Tabu
search \citep{glover1989} mengingat langkah-langkah terakhir dan melarang kembali ke sana, sehingga
tidak berputar-putar di sekitar puncak yang sama.

\section{Bermain melawan lawan}

Dalam permainan dua pemain, keadaan berikutnya juga ditentukan oleh lawan yang berusaha
mengalahkan kita. Algoritme \istilah{minimax} mengasumsikan lawan bermain sempurna: di giliran kita
pilih langkah dengan nilai maksimum, di giliran lawan anggap ia memilih nilai minimum. Contoh dari
\citet{russell2021} memperlihatkannya. Misalkan kita punya tiga langkah, dan lawan menjawab setiap
langkah dengan tiga kemungkinan yang nilainya $(3,12,8)$, $(2,4,6)$, dan $(14,5,2)$. Lawan akan
memilih minimum di setiap cabang, yaitu 3, 2, dan 2, sehingga kita memilih cabang pertama dengan
nilai 3.

\istilah{Alpha-beta pruning} \citep{knuth1975} memperoleh jawaban yang sama tanpa melihat semua daun.
Setelah cabang pertama memberi nilai 3, kita tahu kita bisa mendapat setidaknya 3. Di cabang kedua,
begitu kita melihat daun bernilai 2, kita tahu lawan akan memberi paling banyak 2 di cabang itu,
sehingga sisa daunnya tidak perlu diperiksa. Dengan urutan langkah yang baik, alpha-beta bisa
menggandakan kedalaman pencarian dengan waktu yang sama.

Untuk permainan dengan cabang banyak seperti Go, minimax tidak mungkin dijalankan sampai
akhir permainan. \istilah{Monte Carlo tree search} membangun pohon pencarian secara bertahap dengan
memainkan banyak permainan acak sampai selesai, lalu memusatkan pencarian pada langkah yang
menjanjikan. Aturan UCT \citep{kocsis2006} menyeimbangkan eksploitasi langkah yang sejauh ini
terbaik dengan eksplorasi langkah yang belum banyak dicoba, mengikuti teori bandit. AlphaGo
\citep{silver2016} menggabungkan MCTS dengan jaringan saraf yang menilai posisi dan menyarankan
langkah, lalu AlphaZero \citep{silver2018} mempelajari catur, shogi, dan Go hanya dari bermain
melawan dirinya sendiri. Kuliah menekankan pelajaran yang lebih umum: pembelajaran memberi
intuisi, pencarian memberi ketelitian, dan sistem terkuat memakai keduanya.

\section{Perencanaan}

Perencanaan adalah pencarian dengan bahasa yang lebih kaya. Alih-alih keadaan sebagai kotak hitam,
keadaan dijelaskan dengan fakta-fakta logis, dan setiap aksi punya prasyarat dan efek. Bahasa STRIPS
\citep{fikes1971}, dan penerusnya PDDL, menulis aksi seperti ``pindahkan balok A dari meja ke atas
balok B'' dengan prasyarat bahwa A dan B kosong di atasnya, dan efek bahwa A sekarang di atas B.
Representasi ini memungkinkan heuristik umum yang tidak bergantung pada domain, misalnya dengan
menyelesaikan versi masalah yang dilonggarkan, di mana efek negatif aksi diabaikan. Panjang rencana
untuk masalah yang dilonggarkan memberi perkiraan biaya untuk masalah aslinya.

\section{Pengetahuan dan logika}

Representasi pengetahuan menyimpan fakta tentang dunia dalam bentuk yang bisa dipakai untuk
bernalar. Logika proposisi bekerja dengan pernyataan benar atau salah, sedangkan logika predikat
orde pertama menambahkan objek, relasi, dan kuantor, sehingga bisa menyatakan ``setiap manusia
fana''. Mesin inferensi menurunkan fakta baru dari fakta lama, misalnya dengan resolusi atau
dengan rantai aturan maju dan mundur. Ontologi dan graf pengetahuan menyusun konsep ke dalam
hierarki dan relasi, dan bentuk modernnya dipakai mesin pencari untuk menjawab pertanyaan faktual
secara langsung. Kekuatan pendekatan simbolik adalah ketelitian dan kemampuan menjelaskan
penalarannya. Kelemahannya adalah kerapuhan di hadapan dunia nyata yang berisik dan tidak lengkap,
persis di tempat machine learning unggul.

\section{Keputusan optimal: operations research}

Kuliah operations research membahas keputusan yang bisa dirumuskan sebagai optimasi dengan
kendala. Bentuk paling penting adalah \istilah{linear programming}: maksimalkan fungsi linear dengan
kendala pertidaksamaan linear. Contoh klasik dari \citet{hillier2015} adalah pabrik yang membuat dua
produk dengan keuntungan 3 dan 5 per unit, dan tiga pabrik pendukung dengan kapasitas terbatas:
\begin{equation}
\max\ 3x+5y\quad\text{dengan}\quad x\le4,\quad 2y\le12,\quad 3x+2y\le18,\quad x,y\ge0.
\end{equation}
Daerah yang memenuhi semua kendala adalah sebuah poligon, dan karena fungsi tujuannya linear,
optimumnya selalu berada di salah satu titik sudut. Di sini titik terbaiknya $(x,y)=(2,6)$ dengan
keuntungan 36. Metode simpleks berjalan dari satu titik sudut ke titik sudut tetangga yang lebih
baik sampai tidak ada lagi perbaikan.

Setiap program linear punya pasangan, yaitu program dual, yang variabelnya bisa dibaca sebagai
harga bayangan setiap sumber daya: berapa keuntungan bertambah kalau kapasitas satu pabrik
dinaikkan satu unit. Pada optimum, nilai primal dan dual sama, dan teori dualitas ini menjadi dasar
banyak algoritme optimasi modern.

Kalau variabel harus bilangan bulat, misalnya jumlah truk, masalahnya menjadi \istilah{integer
programming}, yang jauh lebih sulit. \istilah{Branch and bound} memecah masalah menjadi cabang dengan
menambahkan kendala seperti $x\le2$ atau $x\ge3$, memakai solusi relaksasi linear sebagai batas,
dan membuang cabang yang batasnya sudah lebih buruk dari solusi terbaik yang ditemukan.
\istilah{Dynamic programming} \citep{bellman1957} memecah keputusan berurutan menjadi submasalah
dengan prinsip optimalitas: bagian akhir dari kebijakan optimal juga optimal untuk keadaan tempat
bagian itu dimulai. Prinsip yang sama menjadi persamaan Bellman di reinforcement learning
(\cref{ch:rl}). Untuk masalah yang terlalu besar, heuristik seperti simulated annealing dan tabu
search kembali menjadi pilihan, dan analisis keputusan menambahkan ketidakpastian dengan pohon
keputusan dan nilai harapan.

\section{Mengapa semua ini masih penting}

Sistem AI modern jarang murni simbolik atau murni statistik. AlphaZero adalah pencarian pohon yang
dipandu jaringan saraf. Model bahasa yang menyelesaikan soal matematika sering menghasilkan
beberapa kemungkinan langkah lalu memilih yang terbaik, sebuah bentuk pencarian. Robot gudang
memakai A* untuk rute dan pembelajaran untuk persepsi. Dalam simulasi teknik, keputusan tentang
desain dan jadwal operasi sering diselesaikan dengan program linear atau integer yang dibangun di
atas prediksi dari model pengganti. Pembelajaran dan penalaran tidak bersaing. Keduanya saling
melengkapi.

\sumberbab{Bab ini mengikuti kuliah Introduction to AI (sejarah, search, game playing, planning,
representasi pengetahuan, machine learning untuk permainan, filsafat AI) dan kuliah KS Operations
Research (linear programming dan dualitas, integer programming, dynamic programming, heuristik,
analisis keputusan). Buku rujukan utamanya adalah \citet{russell2021} dan \citet{hillier2015}.}
